% capitulos/04-cronograma.tex

O cronograma de execução da pesquisa está organizado em quatro anos,
correspondendo ao prazo máximo estabelecido pelo regulamento do
\abreviatura{PPGCA}{Programa de Pós-Graduação em Ciências Ambientais}
para a obtenção do título de Doutor. As atividades estão distribuídas
conforme a \autoref{tab:cronograma}.

As etapas foram concebidas de forma a respeitar os períodos sazonais
do Cerrado, com as campanhas de campo concentradas nos meses de
transição entre o período seco e o chuvoso (setembro--outubro) e no
início do período chuvoso (novembro--dezembro), quando a umidade do
solo e a atividade biológica se mostram mais representativas das
condições anuais médias \cite{exemplo:artigo:2024}.

\begin{table}[H]
  \centering
  \caption{Cronograma de execução da pesquisa (2023--2027).
           \textbullet\ = atividade prevista; \textbf{\textbullet} = atividade concluída.}
  \label{tab:cronograma}
  \small
  \begin{tabularx}{\textwidth}{Xcccccccccccc}
    \toprule
    \multirow{2}{*}{\textbf{Atividade}}
      & \multicolumn{3}{c}{\textbf{2023}}
      & \multicolumn{3}{c}{\textbf{2024}}
      & \multicolumn{3}{c}{\textbf{2025}}
      & \multicolumn{3}{c}{\textbf{2026/27}} \\
    \cmidrule(lr){2-4}\cmidrule(lr){5-7}\cmidrule(lr){8-10}\cmidrule(lr){11-13}
      & T1 & T2 & T3
      & T1 & T2 & T3
      & T1 & T2 & T3
      & T1 & T2 & T3 \\
    \midrule
    Revisão bibliográfica
      & $\bullet$ & $\bullet$ & $\bullet$
      & $\bullet$ &            &
      &            &            &
      &            &            & \\
    \thinrule
    Qualificação do doutorado
      &            &            & $\bullet$
      &            &            &
      &            &            &
      &            &            & \\
    \thinrule
    Campanha de campo 1 (período seco)
      &            &            & $\bullet$
      &            &            &
      &            &            &
      &            &            & \\
    \thinrule
    Análises laboratoriais -- campanha 1
      &            &            &
      & $\bullet$ & $\bullet$ &
      &            &            &
      &            &            & \\
    \thinrule
    Campanha de campo 2 (transição)
      &            &            &
      &            & $\bullet$ &
      &            &            &
      &            &            & \\
    \thinrule
    Campanha de campo 3 (período chuvoso)
      &            &            &
      &            &            & $\bullet$
      &            &            &
      &            &            & \\
    \thinrule
    Análises laboratoriais -- campanhas 2 e 3
      &            &            &
      &            &            &
      & $\bullet$ & $\bullet$ &
      &            &            & \\
    \thinrule
    Análise estatística e modelagem
      &            &            &
      &            &            &
      &            & $\bullet$ & $\bullet$
      &            &            & \\
    \thinrule
    Redação e submissão de artigos
      &            &            &
      &            &            &
      &            &            & $\bullet$
      & $\bullet$ & $\bullet$ & \\
    \thinrule
    Redação da tese
      &            &            &
      &            &            &
      &            &            &
      & $\bullet$ & $\bullet$ & $\bullet$ \\
    \thinrule
    Defesa pública
      &            &            &
      &            &            &
      &            &            &
      &            &            & $\bullet$ \\
    \bottomrule
  \end{tabularx}
  \fonte{Elaborado pelo autor (2025). T1~=~Jan--Abr; T2~=~Mai--Ago; T3~=~Set--Dez.}
\end{table}

As atividades de revisão bibliográfica e análise de dados são contínuas
ao longo de toda a pesquisa, não se restringindo aos períodos
indicados na tabela. A redação dos manuscritos para publicação ocorre
em paralelo à análise dos dados de cada campanha de campo.
